% $Id: gui_disres.tex 395 2011-11-17 22:34:59Z cphillip $
%
% Written by Jonas Richiardi
%_______________________________________________


\chapter{Results display}
\label{chap:DisRes}
\minitoc

\section{Introduction}

Once a machine (e.g. a classifier or a regression function) has been specified, 
its parameters have been estimated over training data, and its performance has been
evaluated over a test set in cross-validation, it is necessary to examine the outcome
of the whole procedure in details. For example, if the machine cannot perform above
chance, it could be due to an inappropriate experimental paradigm, noisy data, insufficient
amount of data, wrong choice of features, wrong choice of machine. It is important to
recognise that any of these factors could cause the modelling to fail. The results window
can help pinpointing the source of error.

Another important aspect at this stage is to see what the machine has learned - some brain 
areas are probably more informative about class membership than others. For example, in a visual
task, we would expect discriminative information in the occipital lobe. This is 
called \textit{information mapping}, and it is of particular import to be critical at this stage - if
the discriminative weight of a machine is concentrated in the eyes, for example, it is important
to correct the analysis mask that was used to exclude them. The question to ask is
``which voxels drive the modelling, and do they make sense with respect to the experimental
paradigm and neurophysiology'' ?

\section{Launching results display}

Make sure all previous steps have been performed (Data and Design, Chapter~\ref{chap:DataDesign};
Prepare feature set, Chapter~\ref{chap:PrepFeat}; Specify Model and Run Model, Chapter~\ref{chap:ModelCrossV} ).
Optionally, you may want to generate a weight map for your machine (Compute Weights~\ref{chap:ComputeWeights}), 
but this is not mandatory.

In the Review Options panel, press {\tt Display Results} to bring up the results main window.